\section{The Folded-Napkin Return Story}

The modern story usually runs as follows: at a Jewish meal, a master tossed
aside his napkin if finished but folded it to tell a servant that he would
return; Jesus therefore folded the face-cloth to promise his return. Every link
requires proof: that such a custom existed in the relevant Jewish period, that
\emph{soudarion} means the table napkin in that custom, that John's participle
means a deliberate fold, and that the Evangelist expected readers to decode
the gesture.

The exact rabbinic controls inspected for this study do not establish any
link. Mishnah Berakhot 8:3 and Babylonian Talmud Berakhot 52b concern where a
hand-wiping cloth is placed during a meal; Mishnah Sanhedrin 6:1 describes a
cloth waved to halt an execution. None contains a departing master, an
end-of-meal fold, or a promise of return. These three loci are controls, not an
exhaustive search of rabbinic literature. Roman dining evidence likewise
supplies cloths and servants without this code. Patristic and medieval
commentators discuss burial, theft, order, and belief; they do not transmit the
master-servant explanation. The Authorized Version contributes the old English
noun ``napkin'' but says it was ``wrapped together,'' not folded as a signal.

The claim appears in modern preaching and internet circulation, often without
a primary citation. Its neatness helps it travel: one unfamiliar object, one
lost custom, one uplifting message. Those are also warning signs. An explanation
that depends on inaccessible cultural knowledge should become more credible as
its ancient documentation improves. Here the documentation disappears as the
search moves backward.

The Christian hope that Christ will return does not need this support. It is
explicit in Scripture and Creed. Retiring the story protects that hope from a
fabricated historical foundation.

